\documentclass[11pt,a4paper]{article}

% --- Packages ---
\usepackage{lmodern}
\usepackage[T1]{fontenc}
\usepackage[utf8]{inputenc}
\usepackage[margin=28mm]{geometry}
\usepackage{amsmath,amssymb}
\usepackage{booktabs}
\usepackage[colorlinks=true,linkcolor=blue,citecolor=blue,urlcolor=blue]{hyperref}
\usepackage{xcolor}
\usepackage{fancyhdr}
\usepackage{microtype}
\usepackage{enumitem}
\usepackage{titlesec}
\usepackage{graphicx}

% --- Section colours ---
\definecolor{efc}{RGB}{0,80,160}
\titleformat{\section}{\large\bfseries\color{efc}}{}{0em}{}[\titlerule]
\titleformat{\subsection}{\normalsize\bfseries\color{efc}}{}{0em}{}

% --- Header / Footer ---
% Canonical Symbiose Research preprint footer (two-column).
% \PaperDOI is set once and propagated to footer + title block.
% Pre-deposit value: [to be assigned] — search/replace after Figshare DOI returns.
\newcommand{\PaperDOI}{10.6084/m9.figshare.32162706}
\newcommand{\PaperShortTitle}{Sign-flip decomposition}
\pagestyle{fancy}
\fancyhf{}
\renewcommand{\headrulewidth}{0pt}
\fancyfoot[L]{\small Magnusson --- \PaperShortTitle{} \textbullet{} DOI: \PaperDOI}
\fancyfoot[R]{\small \thepage}

% --- Metadata ---
\title{\textbf{Composite-level aggregation in rotation curve model selection:\\
A sign-flip decomposition of $\Delta$AIC for SPARC-175}}

\author{
  Morten Magnusson\\
  \small Symbiose Research, Sandnes, Norway\\
  \small ORCID: \href{https://orcid.org/0009-0002-4860-5095}{0009-0002-4860-5095}\\
  \small DOI: \href{https://doi.org/\PaperDOI}{\PaperDOI}
}
\date{May 2026 \quad|\quad Figshare preprint}

\begin{document}
\maketitle
\thispagestyle{fancy}

% ==============================================================
% Sign convention used throughout this paper:
%   Delta AIC = AIC_NFW - AIC_EFC   (positive = EFC wins)
% This matches Kill-Test v6 (Magnusson 2026b). Earlier drafts of
% this work (v1, v2) used the opposite convention; all numerical
% values in this paper have been transposed to v6 sign.
% ==============================================================

% --- Abstract ---
\begin{quote}
\textbf{Abstract.}
Comparing gravitational models on the SPARC-175 sample using a single
global $\Delta$AIC per galaxy collapses two physically distinct radial
regimes onto one number.  We re-fit all 175 galaxies with the Kill-Test~v6
pipeline, persist per-datapoint model curves, and split each $\Delta$AIC
into an inner ($r<r_c$) and outer ($r\ge r_c$) contribution at
$r_c=3$~kpc.  In $\mathbf{30/171}$ galaxies (\textbf{17.5\%}) the global
sign disagrees with the outer-only sign: a model preferred globally is
disfavoured by the outer points alone.  This sign-flip rate is robust,
varying only between $\mathbf{16.4\%}$ and $\mathbf{22.2\%}$ for
$r_c\in\{2,3,4,5\}$~kpc.  We test whether $\chi^2$-weighting drives the
result: replacing $1/\sigma^2$ weights with uniform weights moves the
sample {\em toward} EFC (from 66.7\% to 86.5\% EFC wins), not away from
it---rejecting the hypothesis that standard weighting is anti-EFC.  The
LATENT regime is the crown witness: 16 galaxies exhibit
$\langle\Delta\text{AIC}_{\text{inner}}\rangle = +153$ (EFC favoured at
small radii) and $\langle\Delta\text{AIC}_{\text{outer}}\rangle = -809$
(NFW overwhelmingly favoured at large radii).  These are not measurement
artefacts; they are a real, regime-stratified failure mode of the
single-statistic comparison.  $\Delta$AIC is a correct statistic at the
wrong aggregation level: a per-regime decomposition is required for
interpretation.  We do not find evidence that standard $\chi^2$ weighting
systematically favours one model over another globally; the effect we
identify is a regime-mixing misclassification at the level of individual
galaxies, not an estimator-level bias against any particular gravitational
model.
\end{quote}

\vspace{4pt}
\noindent\textbf{Keywords:} rotation curves; model selection; AIC decomposition;
SPARC-175; regime stratification; sign-flip diagnostic;
entropy-bounded empiricism

% ==============================================================
\newpage
\section{Introduction}

Rotation-curve model selection compresses each galaxy to a small set of
summary statistics: $\chi^2$, AIC, BIC.  Comparative analyses across the
SPARC sample (Lelli et~al.\ 2016) routinely report a per-galaxy
$\Delta$AIC and aggregate over the 175-galaxy sample to declare a winner.
This paper asks a methodological question: does global $\Delta$AIC
preserve the radial information that distinguishes inner-region
(baryon-dominated) from outer-region (gravity-dominated) physics, or
does aggregation collapse two regimes into a single number that no
single physical mechanism can satisfy?

Our previous draft (v2; Magnusson 2026e, DOI:
\href{https://doi.org/10.6084/m9.figshare.32019723}{10.6084/m9.figshare.32019723})
framed the answer as a measurement-level bias: inner-region weighting
could be argued to favour EFC over NFW.  The
analysis below rejects that framing.  When weights are removed, the
sample moves further toward EFC, not away from it.  The real signal is
not in {\em how} points are weighted; it is in {\em where on the
rotation curve} the model wins or loses.  In 17.5\% of SPARC galaxies
the inner and outer regions disagree on which model is preferred,
producing a sign-flip in $\Delta$AIC that no global statistic can detect.

We adopt throughout the Kill-Test v6 sign convention
\begin{equation}
  \Delta\text{AIC} \equiv \text{AIC}_{\text{NFW}} - \text{AIC}_{\text{EFC}},
  \qquad \text{positive}=\text{EFC wins.}
  \label{eq:convention}
\end{equation}
This is the inverse of the convention used in earlier drafts of this work.
All numerical values reported here have been transposed.

The rest of the paper proceeds as follows.  Section~\ref{sec:methods}
specifies the fits, the entropy-bounded empiricism (EBE) regime
classification used as scaffolding, and the inner/outer decomposition.
Section~\ref{sec:signflip} reports the headline sign-flip statistic and
its $r_c$-robustness.  Section~\ref{sec:reweight} reports the
reweighting test honestly, including the result that contradicts the v2
narrative.  Section~\ref{sec:regime} stratifies by regime; the LATENT
sub-sample (\S\ref{sec:latent}) is presented as the crown witness for
the aggregation argument.  Section~\ref{sec:discussion} discusses
implications for $\Delta$AIC reporting practice and the proper level of
empirical claim that the SPARC sample supports.

% ==============================================================
\section{Methods}
\label{sec:methods}

\subsection{Kill-Test v6 fits}

For each SPARC galaxy with $\ge 5$ valid data points, we fit two
rotation curve models by differential evolution
($\text{seed}=42$, Sobol initialisation, $\text{maxiter}=500$,
$\text{tol}=10^{-8}$):
\begin{align}
  V_{\text{EFC}}(r) &= v_{\text{flat}}\sqrt{\,1-\exp[-(r/r_{\text{turnon}})^{p}]\,},
  & k_{\text{EFC}} &= 3, \\
  V_{\text{NFW}}(r) &= V_{200}\sqrt{[\ln(1+s)-s/(1+s)]\,/\,
                        [\ln(1+c)-c/(1+c)]\,/\,(r/r_{200})}, &  k_{\text{NFW}} &= 2,
\end{align}
where $s=cr/r_{200}$ and $r_{200}=V_{200}/(10\,H_0\!\cdot\!10^{-3})$
with $H_0=67.4$~km\,s$^{-1}$\,Mpc$^{-1}$.  Bounds and AIC computation
match Kill-Test v6 verbatim.  Of 175 galaxies, 171 converged for both
models and are retained for analysis.  The per-datapoint model
velocities $V_{\text{EFC}}(r_i)$ and $V_{\text{NFW}}(r_i)$ at the
best-fit parameters are persisted via a \texttt{-{}-save-curves} flag
(see \S\ref{sec:repro}).

\subsection{Regime classification (EBE scaffolding)}

We retain the entropy-bounded empiricism (EBE; Magnusson 2026d) regime
labels FLOW, TRANSITION, LATENT, used here as a stratification scaffold,
not a physical prediction:
\begin{equation}
  \text{regime} =
  \begin{cases}
    \text{FLOW} & \text{if } \Delta\text{AIC}\ge 10 \text{ and } \chi^2_r({\text{EFC}}) < 20 \\
    \text{LATENT} & \text{if } \Delta\text{AIC}\le -10 \text{ or } \chi^2_r({\text{EFC}}) > 50 \\
    \text{TRANSITION} & \text{otherwise.}
  \end{cases}
\end{equation}
EBE asserts that empirical claims should not aggregate across regimes
that bound observation differently; the present paper is a test of
whether ignoring that asymmetry materially distorts SPARC-level
conclusions.  EBE is the framework, not the headline.

\subsection{Inner/outer decomposition}
\label{sec:decomp}

We define a critical radius $r_c$ separating inner from outer points and
the corresponding regime contributions
\begin{align}
  \Delta\text{AIC}_{\text{inner}}(r_c) &=
    \sum_{r_i<r_c} w_i\bigl[(V_{\text{obs},i}-V_{\text{NFW},i})^2 -
                              (V_{\text{obs},i}-V_{\text{EFC},i})^2\bigr]
    + 2(k_{\text{NFW}}-k_{\text{EFC}}), \\
  \Delta\text{AIC}_{\text{outer}}(r_c) &=
    \sum_{r_i\ge r_c} w_i\bigl[\cdots\bigr] + 2(k_{\text{NFW}}-k_{\text{EFC}}),
\end{align}
with $w_i=1/\sigma_i^2$ unless stated otherwise.  The $2\Delta k$ penalty
is applied to each split rather than refitting on subsets; this is a
stratified contribution analysis, not a re-fit.  We define the global
quantity by the same expression with the radial mask removed; by
construction
$\Delta\text{AIC}_{\text{global}}(r_c) = \Delta\text{AIC}_{\text{inner}}(r_c)
                                          + \Delta\text{AIC}_{\text{outer}}(r_c)
                                          - 2(k_{\text{NFW}}-k_{\text{EFC}})$.
The recomputed global value differs from the Kill-Test v6 number only by
floating-point round-off; we use the recomputed value throughout for
internal consistency.

\paragraph{Heuristic nature of the split.}
The decomposition above evaluates contributions to the $\chi^2$ sum at
the {\em globally} best-fit model parameters rather than refitting each
model independently on the radial subsets $r<r_c$ and $r\ge r_c$.  This
is a deliberate choice: the diagnostic question is whether the
single-number $\Delta$AIC summary preserves the radial information that
the model-on-full-curve fit had access to, and that question is answered
exactly by stratifying the same $\chi^2$ sum by radius.  Independent
subset refits would inflate parameter uncertainty (especially for $n_{\text{outer}}<10$),
introduce degeneracies between $V_{200}$ and $c$ in the truncated NFW
fit, and would no longer be a decomposition of the original AIC
comparison but a different comparison altogether.  We treat the split
therefore as a heuristic AIC contribution analysis valid for the
diagnostic purpose stated in this paper; subset-level refits, with the
parameter-uncertainty discussion they require, are reserved for
follow-up work and are not a precondition for the conclusions reported
here.

\subsection{Sign-flip diagnostic}

A galaxy exhibits a sign-flip at $r_c$ when
\begin{equation}
  \mathrm{sign}\bigl[\Delta\text{AIC}_{\text{global}}\bigr]
  \ne
  \mathrm{sign}\bigl[\Delta\text{AIC}_{\text{outer}}\bigr],
\end{equation}
i.e.\ the model preferred by the full curve is disfavoured by the outer
points alone.  This is a direct, model-agnostic test of whether
single-number $\Delta$AIC has thrown away information that the data
contained.

% ==============================================================
\section{Sign-flip detection}
\label{sec:signflip}

At $r_c=3$~kpc, $\mathbf{30/171}$ galaxies (\textbf{17.5\%}) satisfy the
sign-flip condition.  Figure~\ref{fig:signflip} plots
$\Delta\text{AIC}_{\text{outer}}$ against $\Delta\text{AIC}_{\text{global}}$,
coloured by EBE regime.  Sign-flip galaxies populate the off-quadrant
diagonals.

\begin{figure}[h]
\centering
\includegraphics[width=0.85\textwidth]{rcmp_run/figs/rcmp_signflip.png}
\caption{Outer ($r\ge 3$~kpc) versus global $\Delta$AIC for 171
SPARC galaxies.  Convention: positive = EFC wins.  Identity line
(no decomposition effect) shown grey; off-diagonal points are
where outer and global disagree.  Galaxies in the upper-left and
lower-right quadrants are the 30 sign-flips (17.5\%).  The lower
left is the LATENT regime, where global and outer agree but inner
disagrees (\S\ref{sec:latent}).}
\label{fig:signflip}
\end{figure}

The Spearman correlation between $\Delta\text{AIC}_{\text{global}}$ and
$\Delta\text{AIC}_{\text{outer}}$ is $\rho=+0.831$ ($p=2.4\times10^{-42}$),
confirming that the outer signal does drive most of the global statistic
{\em on average}, while leaving room for the 17.5\% population in which
they disagree.

\subsection{Robustness to $r_c$}

\begin{table}[h]
\centering
\small
\begin{tabular}{cccc}
\toprule
$r_c$ (kpc) & $n_{\text{sign-flip}}$ & fraction & $n$ with both inner/outer points \\
\midrule
2.0 & 28 & 16.4\% & 154 \\
3.0 & 30 & 17.5\% & 161 \\
4.0 & 35 & 20.5\% & 167 \\
5.0 & 38 & 22.2\% & 167 \\
\bottomrule
\end{tabular}
\caption{Sign-flip count is stable to within a few percentage points
across the range of plausible $r_c$ values.  The result is not a
knife-edge of one chosen radius.}
\label{tab:rc}
\end{table}

% ==============================================================
\section{The reweighting test (and its honest result)}
\label{sec:reweight}

To separate aggregation effects from weighting effects, we recompute
the global $\Delta\text{AIC}$ under three weight schemes:
\begin{description}[leftmargin=2em]
  \item[Standard] $w_i=1/\sigma_i^2$, identical to the Kill-Test v6 fit.
  \item[Uniform]  $w_i=1$, every datapoint contributes equally.
  \item[Capped]   $w_i=\min(1/\sigma_i^2,\,\text{med}(1/\sigma_j^2))$,
                  capping the most precise points at the median weight.
\end{description}

\begin{table}[h]
\centering
\small
\begin{tabular}{lrrr}
\toprule
weight scheme & mean $\Delta\text{AIC}$ & median & EFC wins ($\Delta\text{AIC}>0$) \\
\midrule
standard ($1/\sigma^2$) & $+40.4$  & $+6.21$  & 114/171 (66.7\%) \\
uniform                 & $+4416$  & $+148.9$ & \textbf{148/171} (\textbf{86.5\%}) \\
capped (median cap)     & $+36.3$  & $+4.47$  & 110/171 (64.3\%) \\
\bottomrule
\end{tabular}
\caption{Replacing $1/\sigma^2$ weights with uniform weights does not
attenuate the EFC advantage; it amplifies it.  This rejects the v2
hypothesis that standard $\chi^2$ weighting is anti-EFC.  Capped weights
are within 2 percentage points of standard, ruling out a few
high-precision points as the driver.}
\label{tab:reweight}
\end{table}

The honest interpretation is that weighting choice {\em does} swing the
sample by 20 percentage points, but in the opposite direction from what
the v2 narrative claimed.  Standard $\chi^2$ weighting is, if anything,
conservatively pro-NFW.  The reweighting test settles a dispute about
weight, but it does not explain the sign-flips: those persist under all
three schemes (the table reports global only; the per-galaxy
inner/outer split shows the same population of sign-flippers under
uniform weighting, with magnitudes inflated).

% ==============================================================
\section{Regime stratification}
\label{sec:regime}

\begin{table}[h]
\centering
\small
\begin{tabular}{lrrrrr}
\toprule
regime & $n$ & $\langle\Delta\text{AIC}_{\text{global}}\rangle$ & $\langle\Delta\text{AIC}_{\text{inner}}\rangle$ & $\langle\Delta\text{AIC}_{\text{outer}}\rangle$ & sign-flips \\
\midrule
FLOW       & 71 & $+231$ & $+189$ & $+50$  & 10 / 71 (14\%) \\
TRANSITION & 84 & $+20$  & $+14$  & $+6$   & 20 / 84 (24\%) \\
LATENT     & 16 & $-702$ & $+153$ & $-809$ & 0 / 16 (0\%) \\
\bottomrule
\end{tabular}
\caption{Regime-stratified decomposition at $r_c=3$~kpc.  All entries in
Kill-Test v6 sign convention (positive = EFC wins).  Sign-flip
populations are concentrated in TRANSITION (24\%) where global is near
zero, and absent in LATENT where the outer term is dominant by orders
of magnitude.}
\label{tab:regime}
\end{table}

The pattern is regime-specific.  TRANSITION galaxies have global
$\Delta$AIC clustered near zero; the inner/outer split has room to
disagree, and 24\% do.  FLOW galaxies are EFC-favoured in both halves
($\langle\Delta\text{AIC}_{\text{inner}}\rangle=+189$,
$\langle\Delta\text{AIC}_{\text{outer}}\rangle=+50$), but a tail of 10
galaxies have outer points that prefer NFW even when the global
statistic is firmly with EFC---these are the high-magnitude sign-flips
of Table~\ref{tab:topflips}.  LATENT galaxies do not flip, because the
outer term is so large in magnitude that it sets the sign of the
global; what they do show, instead, is the most informative feature
of the entire dataset.

\begin{table}[h]
\centering
\small
\begin{tabular}{lrrrrr}
\toprule
galaxy & regime & global & inner & outer \\
\midrule
UGC03546 & FLOW & $+90.7$  & $+91.8$  & $-3.1$  \\
NGC7814  & FLOW & $+85.8$  & $+97.0$  & $-13.2$ \\
KK98-251 & FLOW & $+41.0$  & $+41.3$  & $-2.3$  \\
NGC6946  & FLOW & $+34.4$  & $+127.1$ & $-94.6$ \\
UGC08286 & FLOW & $+27.7$  & $+33.7$  & $-8.1$  \\
UGC04305 & FLOW & $+26.7$  & $+26.7$  & $-2.0$  \\
UGC07151 & FLOW & $+23.9$  & $+28.0$  & $-6.1$  \\
UGC07603 & FLOW & $+14.0$  & $+13.0$  & $-0.9$  \\
UGC00191 & FLOW & $+12.5$  & $+11.5$  & $-1.0$  \\
NGC4157  & FLOW & $+10.9$  & $+9.1$   & $-0.2$  \\
\bottomrule
\end{tabular}
\caption{Top ten sign-flip galaxies by $|\Delta\text{AIC}_{\text{global}}|$.
All are FLOW.  Each shows a global EFC win driven entirely or
predominantly by inner-region structure; the outer points alone prefer
NFW.  NGC6946 is the cleanest example: $\Delta\text{AIC}_{\text{inner}}=+127$
(EFC) versus $\Delta\text{AIC}_{\text{outer}}=-95$ (NFW).}
\label{tab:topflips}
\end{table}

% ==============================================================
\section{The LATENT regime as crown witness}
\label{sec:latent}

LATENT contains 16 galaxies whose global $\Delta\text{AIC}$ averages
$-702$: NFW wins by a large margin.  The decomposition tells a sharper
story.  At $r<3$~kpc, the mean inner $\Delta\text{AIC}$ is $\mathbf{+153}$:
EFC fits the inner regions of LATENT galaxies at the same level of
preference seen in FLOW.  At $r\ge3$~kpc, the mean outer $\Delta\text{AIC}$
is $\mathbf{-809}$: NFW wins by roughly four orders of magnitude in the
outer regions.

This is not a measurement artefact.  The errors in LATENT galaxies are
not pathological, the inner-fit residuals are not anomalous, and the
result is independent of $r_c$ within the range tested.  It is a real
physical statement: the EFC functional form, with its three free
parameters, fits the inner kinematics of LATENT galaxies as well as it
fits any FLOW galaxy, but it cannot accommodate the outer rotation
behaviour of these systems.  NFW, conversely, fits the outer regions
with overwhelming evidence and is mediocre on the inner.

The single-statistic $\Delta\text{AIC}$ erases this distinction.  A
LATENT galaxy with $\Delta\text{AIC}_{\text{global}}=-700$ is reported
as ``NFW preferred,'' suggesting the outer NFW success extends to all
radii.  It does not.  The same galaxy is also a galaxy whose inner
3~kpc strongly disagrees, and that disagreement carries empirical content
of its own.  Reporting only the global suppresses 50\% of the radial
information; reporting only the regime label, in turn, suppresses the
sign of the inner term.  Both are needed.

This is also the locus where EBE (Magnusson 2026d) earns its scaffolding
role: LATENT and FLOW are distinct empirical regimes, and aggregating
them into a sample mean is a category error precisely of the kind EBE
formalises.  The LATENT result is the strongest single piece of evidence
in this paper that single-number $\Delta\text{AIC}$ aggregation is the
wrong empirical level for SPARC-style comparison.

% ==============================================================
\section{Discussion}
\label{sec:discussion}

\subsection{$\Delta$AIC is a correct statistic at the wrong aggregation level}

Nothing in this analysis impugns AIC itself.  $\Delta\text{AIC}$ between
two models on a given sum-of-squared-residuals likelihood remains a
valid, well-grounded comparison.  The result of this paper is purely
about the {\em level} at which the comparison is reported.  When 17.5\%
of galaxies show inner/outer disagreement and a regime-defined
sub-sample shows order-of-magnitude divergence between halves, the
single-number-per-galaxy summary destroys information that the data
contained.  The fix is not a different statistic; it is a different
aggregation.  We recommend reporting
$(\Delta\text{AIC}_{\text{inner}}, \Delta\text{AIC}_{\text{outer}})$
as a pair, with $r_c$ stated explicitly, alongside the global value.

\subsection{What is rejected, what is not}

We explicitly reject the v2 framing of this work that standard
$\chi^2$ weighting is anti-EFC.  Uniform weighting moves the sample
{\em toward} EFC by 20 percentage points (Table~\ref{tab:reweight}).
The composite-bias narrative as a measurement-level critique of NFW does
not survive.

We do {\em not} retract the headline of this paper, only its frame.  The
finding---that single-number $\Delta\text{AIC}$ over the SPARC sample
collapses two regimes onto one summary, in 17.5\% of galaxies producing
a sign-flip---is methodologically robust and independent of any
particular gravitational model.  It applies in equal force to MOND/RAR
comparisons (Li et~al.\ 2020; McGaugh 2020), to NFW vs.\ alternative-DM
comparisons (Oman et~al.\ 2015), and to any future model comparison on
SPARC-style rotation curves.

\subsection{The LATENT result is physical}

We make no claim that the LATENT outer collapse of EFC is an artefact of
weighting, error inflation, or fit pathology.  In LATENT galaxies, EFC
{\em as currently formulated} fits the inner curve and fails the outer.
This is a real boundary on the EFC functional form, and we report it as
such.  The inner success is equally real, and equally informative for
future model development.

\subsection{Limitations}

(i) Heuristic AIC contribution analysis; subset-refit deferred to
follow-up work (rationale and trade-offs in \S\ref{sec:decomp}).
(ii) $r_c=3$~kpc is chosen for visual separability of regimes; we
report the $r_c\in\{2,4,5\}$ robustness in Table~\ref{tab:rc}.
(iii) The 4 galaxies that did not converge in the Kill-Test pipeline
are excluded; these are short rotation curves with $<5$ valid points
and are unlikely to skew the aggregate.

% ==============================================================
\section{Reproducibility}
\label{sec:repro}

All analysis code and intermediate data are deposited with this preprint:

\begin{itemize}[leftmargin=2em,itemsep=2pt]
  \item \texttt{sparc175\_killtest\_universality\_curves.py} ---
    Kill-Test v6 fits with the \texttt{-{}-save-curves} extension
    (per-datapoint model curves persisted to JSON).
  \item \texttt{rcmp\_decomposition.py} ---
    inner/outer split, sign-flip detection, reweighting test,
    $r_c$-robustness scan.
  \item \texttt{rcmp\_run/sparc175\_model\_curves.json} ---
    per-galaxy $V_{\text{EFC}}(r_i)$, $V_{\text{NFW}}(r_i)$,
    $V_{\text{obs}}(r_i)$, $\sigma_i$, and baryonic components
    $V_{\text{gas}}, V_{\text{disk}}, V_{\text{bulge}}$.
  \item \texttt{rcmp\_run/rcmp\_decomposition.json} ---
    per-galaxy
    $(\Delta\text{AIC}_{\text{global}},\Delta\text{AIC}_{\text{inner}},\Delta\text{AIC}_{\text{outer}})$
    plus reweighting variants.
  \item \texttt{rcmp\_run/rcmp\_summary.json},
        \texttt{rcmp\_signflip\_table.json},
        \texttt{rcmp\_robustness\_rc.json} --- aggregates.
\end{itemize}
Input data: SPARC rotation curves (Lelli et~al.\ 2016).  Frozen
parameters: $H_0=67.4$~km\,s$^{-1}$\,Mpc$^{-1}$, $r_c=3$~kpc,
$\sigma_{\min}=1.0$~km\,s$^{-1}$, EFC bounds and NFW concentration prior
identical to Kill-Test v6 (Magnusson 2026b).

% ==============================================================
\section*{References}

\sloppy
\begin{small}
\begin{enumerate}[leftmargin=2em,label={[\arabic*]},itemsep=2pt]

\item Lelli, F., McGaugh, S.~S., \& Schombert, J.~M.\ (2016).
  SPARC: Mass models for 175 disk galaxies with \textit{Spitzer} photometry
  and accurate rotation curves.
  \textit{AJ}, 152, 157.
  \href{https://doi.org/10.3847/0004-6256/152/6/157}{doi:10.3847/0004-6256/152/6/157}

\item Oman, K.~A.\ et~al.\ (2015).
  The unexpected diversity of dwarf galaxy rotation curves.
  \textit{MNRAS}, 452, 3650.
  \href{https://doi.org/10.1093/mnras/stv1504}{doi:10.1093/mnras/stv1504}

\item Li, P.\ et~al.\ (2020).
  Fitting the radial acceleration relation to individual SPARC galaxies.
  \textit{A\&A}, 634, A71.
  \href{https://doi.org/10.1051/0004-6361/201936368}{doi:10.1051/0004-6361/201936368}

\item McGaugh, S.~S.\ (2020).
  The radial acceleration relation in galaxies and its implications.
  \textit{Galaxies}, 8, 35.

\item Magnusson, M.\ (2026a).
  Comprehensive analysis of 175 SPARC galaxies demonstrating
  regime-dependent validity in rotation curve modelling.
  Figshare.
  \href{https://doi.org/10.6084/m9.figshare.31045126}{doi:10.6084/m9.figshare.31045126}

\item Magnusson, M.\ (2026b).
  Kill-Test v6 Universality: EFC vs NFW across 175 SPARC galaxies.
  Figshare.
  \href{https://doi.org/10.6084/m9.figshare.31986762}{doi:10.6084/m9.figshare.31986762}

\item Magnusson, M.\ (2026c).
  ISW Void Sign-Flip Observational Pipeline From EFC Theory to Data
  Confrontation.
  Figshare.
  \href{https://doi.org/10.6084/m9.figshare.31990212}{doi:10.6084/m9.figshare.31990212}

\item Magnusson, M.\ (2026d).
  Entropy-Bounded Empiricism (EBE): SPARC-175 Complete Documentation.
  Figshare.
  \href{https://doi.org/10.6084/m9.figshare.31047703}{doi:10.6084/m9.figshare.31047703}

\item Magnusson, M.\ (2026e).
  Component-weighted bias in rotation curve model selection: A
  proxy-chain audit of $\chi^2$-based $\Delta$AIC statistics
  [retracted by present work; see \S\ref{sec:reweight} and \S\ref{sec:discussion}].
  Figshare.
  \href{https://doi.org/10.6084/m9.figshare.32019723}{doi:10.6084/m9.figshare.32019723}

\end{enumerate}
\end{small}

\end{document}
